\documentclass{article}
\usepackage[UTF8]{ctex}
\usepackage{amsmath,amssymb,amsfonts}
\usepackage{enumitem}
\usepackage{geometry}
\geometry{a4paper, margin=1in}
% === 错题本排版优化（不参与题目标识） ===
\setlength{\parindent}{0pt}
\setlength{\parskip}{0.35em plus 0.10em minus 0.08em}
\setlength{\emergencystretch}{3em}
\setlist{itemsep=0.25em, topsep=0.40em, parsep=0pt}
\allowdisplaybreaks[2]
% === 错题本排版优化结束 ===
\begin{document}
\section*{电子科技大学2026年数学分析真题}
\section*{169．电子科技大学2026年数学分析试卷}
\begin{itemize}

\subsection{题目1}
  \item[1．] 设 \(a>0\) ，求 \(\lim _{n \rightarrow \infty} n(\sqrt[n]{a}-\sqrt[2 n]{a})=\) \(\_\_\_\_\)．

\subsection{题目2}
  \item[2．] 求 \(\lim _{x \rightarrow+\infty} \frac{1}{x} \int_{\frac{1}{x}}^{1} \frac{\cos 2 t}{t^{2}} \mathrm{~d} t=\) \(\_\_\_\_\)．
\end{itemize}

\subsection{题目3}
3 ．求 \(\int \sqrt{x} e^{\sqrt{x}} \mathrm{~d} x=\) \(\_\_\_\_\)．

\subsection{题目4}
4．令 \(u=\ln (x+y)\) ，则 \(\left.\mathrm{d} u\right|_{(1,2)}=\) \(\_\_\_\_\)．

\subsection{题目5}
5．求 \(\sum_{n=1}^{\infty} \frac{n!}{n^{n}} x^{n}\) 的收敛半径 \(\_\_\_\_\)．

\subsection{题目6}
6．求 \(\int_{0}^{+\infty} \frac{\cos x-\cos 2 x}{x^{2}} \mathrm{~d} x=\) \(\_\_\_\_\)．

\subsection{题目7}
7．已知 \(f(x)\) 在 \([0,+\infty)\) 上连续，且 \(\lim _{x \rightarrow+\infty} f(x)=A\) ，求 \(\lim _{n \rightarrow \infty} \int_{0}^{1} f(n x) \mathrm{d} x\) ．

\subsection{题目8}
8．设 \(f(t)=\iint_{0 \leq x \leq t, 0 \leq y \leq t} e^{-\frac{t x}{v^{2}}} \mathrm{~d} x \mathrm{~d} y, t>0\) ，求 \(f^{\prime}(t)\) ．

\subsection{题目9}
9．已知曲面

\[
\left\{\begin{array}{l}
x=\frac{1}{2}\left(e^{u}+e^{-u}\right) \cos v \\
y=\frac{1}{2}\left(e^{u}+e^{-u}\right) \sin v \\
z=\frac{1}{2}\left(e^{u}-e^{-u}\right)
\end{array}\right.
\]

求曲面在点 \((u, v)=\left(0, \frac{\pi}{4}\right)\) 处的切平面和法线方程．

\subsection{题目10}
10．求曲面积分

\[
\iint_{\sum} \frac{1}{x} \mathrm{~d} y \mathrm{~d} z+\frac{1}{y} \mathrm{~d} z \mathrm{~d} x+\frac{1}{z} \mathrm{~d} x \mathrm{~d} y .
\]

其中 \(\sum\) 为椭球面 \(\frac{x^{2}}{a^{2}}+\frac{y^{2}}{b^{2}}+\frac{z^{2}}{c^{2}}=1\) ，方向取外侧。

\subsection{题目11}
11．设 \(\mathbb{R}\) 为实数集，\(F \subset \mathbb{R}\) 为非空闭集，定义 \(d(x)=\inf _{t \in F}|x-t|\) ，证明 \(d(x)\) 为连续函数．

\subsection{题目12}
12．设 \(f(x) \in R[0,1]\) ，且 \(f\) 在 \(x=0\) 处连续，证明： \(\lim _{n \rightarrow \infty} \int_{0}^{1} f\left(x^{n}\right) \mathrm{d} x=f(0)\) ．

\subsection{题目13}
13．\(x \in[-\pi, \pi]\) ，判断 \(\sum_{n=2}^{\infty} \frac{\cos n x}{\ln n}\) 的敛散性（绝对，条件，发散），并证明其不是某 Riemann 可积函数在\([-\pi, \pi]\) 上的 Fourier 级数．

\subsection{题目14}
14．设 \(K\) 是一个非空紧集，\(\left\{f_{n}(x)\right\}\) 为 \(K\) 上的连续函数列，且对任意的 \(x \in K,\left\{f_{n}(x)\right\}\) 关于 \(n\) 单调递减趋于 0 ，证明：\(\left\{f_{n}(x)\right\}\) 在 \(K\) 上一致收敛于 \(f(x)=0\) ．

\subsection{题目15}
15．设 \(\Omega\) 为 \(\mathbb{R}^{2}\) 上的光滑简单连通区域，\(f(\mathbf{x}) \in C^{2}(\Omega)\) ，若 \(\frac{\partial^{2} f}{\partial x_{1}^{2}}+\frac{\partial^{2} f}{\partial x_{2}^{2}}=0, \forall \mathbf{x}=\left(x_{1}, x_{2}\right) \in \Omega\) ，则称 \(f\)为 \(\Omega\) 上的调和函数。\\
（1）证明： \(\iint_{\Omega}\left(\frac{\partial^{2} f}{\partial x_{1}^{2}}+\frac{\partial^{2} f}{\partial x_{2}^{2}}\right) \mathrm{d} \sigma=\int_{\partial \Omega} \frac{\partial f}{\partial n} \mathrm{~d} s\) ，其中 \(n\) 是 \(\partial \Omega\) 的单位外法向量．\\
（2）证明：\(f\) 为 \(\Omega\) 上的调和函数当且仅当

\[
f(x)=\frac{1}{2 \pi r} \int_{\partial B(x, r)} f(x) \mathrm{d} s=\frac{1}{\pi r^{2}} \iint_{B(x, r)} f(x) \mathrm{d} \sigma
\]

其中 \(B(x, r)\) 是以 \(x\) 为圆心以 \(r\) 为半径的闭球．
\end{document}
